% ============ 第3章 §3.2 Hardy--Littlewood 极大函数的 $L^{p}$ 有界性 (原书 pp.67-70) ============

\section{\texorpdfstring{Hardy--Littlewood 极大函数的 $L^{p}$ 有界性}{Hardy-Littlewood 极大函数的 Lp 有界性}}\label{sec:3.2}
\index{H@Hardy--Littlewood 的 $L^{p}$ 不等式}


对 $1 \leq p < \infty$, (可测)函数 $f$ 的 $L^{p}$ 范数是由
\[
\|f\|_{L^{p}} = \left(\int_{\mathbb{R}^{n}}|f(x)|^{p}\dd x\right)^{1/p}
\]
给出的. 一般地, 对算子 $S$, 说 $S$ 是 $L^{p}$ 有界的, 是指存在常数 $C > 0$, 使得对一切 $f \in L^{p}(\mathbb{R}^{n})$, 不等式 $\|Sf\|_{L^{p}} \leq C\|f\|_{L^{p}}$ 成立. 我们已经知道 Hardy--Littlewood 极大算子不是 $L^{1}$ 有界的, 且是 $L^{\infty}$ 有界的. Hardy--Littlewood 极大算子不是 $L^{1}$ 有界的, 但取而代之, 有称为弱 (1,1) 不等式的不等式成立. 这里, 证明当 $1 < p < \infty$ 时, Hardy--Littlewood 极大算子是 $L^{p}$ 有界的.

从上面的定义式可以看出, $L^{p}$ 范数是难处理的. 这是因为会出现 $p$ 分之一次方与 $p$ 次方. 在 $p = 2$ 的情形, 由于 $|f(x)|^{2} = f(x)\overline{f(x)}$, 或许容易处理, 但除去这种特殊情形, 一般地要对这个量作代数的解释是困难的. 因此, 首先说明不借助函数的 $p$ 次方来表出 $L^{p}$ 范数的方法, 并以此为基础证明 Hardy--Littlewood 极大算子是 $L^{p}$ 有界的.

将定理~\ref{thm:1.3} 改写到函数的情形, 就得到下面的定理.

\begin{theorem}\label{thm:3.3}
设 $1 \leq p < \infty$. 此时, 对可测函数 $f: \mathbb{R}^{n} \to \mathbb{C}$, 有
\begin{equation}\tag{3.8}
\|f\|_{L^{p}} = \left(\int_{0}^{\infty}p\lambda^{p-1}|\{x \in \mathbb{R}^{n}: |f(x)| > \lambda\}|\dd\lambda\right)^{1/p}
\label{eq:3.8}
\end{equation}
成立.
\end{theorem}

\begin{proof}
对非负值函数使用 Fubini 定理, 直接计算.
\begin{align*}
\int_{\mathbb{R}^{n}}|f(x)|^{p}\dd x &= \int_{\mathbb{R}^{n}}\left(\int_{0}^{|f(x)|}p\lambda^{p-1}\dd\lambda\right)\dd x \\
&= \int_{\mathbb{R}^{n}}\left(\int_{0}^{\infty}p\lambda^{p-1}\chi_{\{(x,\lambda): |f(x)| > \lambda\}}(x,\lambda)\dd\lambda\right)\dd x \\
&= \int_{0}^{\infty}\left(\int_{\mathbb{R}^{n}}p\lambda^{p-1}\chi_{\{(x,\lambda): |f(x)| > \lambda\}}(x,\lambda)\dd x\right)\dd\lambda \\
&= \int_{0}^{\infty}p\lambda^{p-1}\left(\int_{\mathbb{R}^{n}}\chi_{\{(x,\lambda): |f(x)| > \lambda\}}(x,\lambda)\dd x\right)\dd\lambda \\
&= \int_{0}^{\infty}p\lambda^{p-1}|\{x \in \mathbb{R}^{n}: |f(x)| > \lambda\}|\dd\lambda.
\end{align*}
对此等式两边取 $p$ 次方根, 即得式~\eqref{eq:3.8}.
\end{proof}

以定理~\ref{thm:3.3} 为基础考察 $M$ 的 $L^{p}$ 有界性. 使用前一节所证明的
\[
|\{Mf > \lambda\}| \leq \frac{3^{n}}{\lambda}\int_{\mathbb{R}^{n}}|f(x)|\dd x, \quad Mf(x) \leq \|f\|_{L^{\infty}}.
\]

\begin{theorem}[Hardy--Littlewood $L^{p}$ 不等式]\label{thm:3.4}
设 $1 < p < \infty$. 此时,
\begin{equation}\tag{3.9}
\|Mf\|_{L^{p}} \leq \left(\frac{p2^{p}\cdot 3^{n}}{p-1}\right)^{1/p}\|f\|_{L^{p}}, \quad f \in L^{p}(\mathbb{R}^{n})
\label{eq:3.9}
\end{equation}
成立.
\end{theorem}

这里出现了
\[
\left(\frac{p2^{p}\cdot 3^{n}}{p-1}\right)^{1/p}
\]
这样的数, 但它无需记忆. 确实存在需要这种常数信息的局面, 但在本书以及相当一般的调和分析的讨论中, 很多情况下没有必要追踪常数. 为什么不必记忆常数呢, 是因为存在定理~\ref{thm:2.22}.

\begin{proof}
首先, 如叙述这个定理之前所述, 注意对一般的可测函数 $h$ 与 $\lambda > 0$, 当 $\|h\|_{\infty} \leq \lambda$ 时, 函数不等式 $Mh \leq \lambda$ 成立. 由这一点与次可加性 $M[g+h] \leq Mg + Mh$, 有 $\{Mf > 2\lambda\} \subset \{M[\chi_{\{|f| > \lambda\}}f] > \lambda\}$. 实际上, 令 $g = \chi_{\{|f| > \lambda\}}f$, $h = f - g = \chi_{\{|f| \leq \lambda\}}f$, 则 $Mh \leq \lambda$. 因此, 若 $x \in \mathbb{R}^{n}$ 满足 $Mf(x) > 2\lambda$, 则
\[
2\lambda < Mf(x) \leq M[\chi_{\{|f| > \lambda\}}f](x) + Mh(x) \leq M[\chi_{\{|f| > \lambda\}}f](x) + \lambda
\]
成立. 所以得到 $\lambda < M[\chi_{\{|f| > \lambda\}}f](x)$. 以上确证了包含关系
\[
\{Mf > 2\lambda\} \subset \{M[\chi_{\{|f| > \lambda\}}f] > \lambda\}.
\]

由这个包含关系与弱 (1,1) 有界性,
\begin{equation}\tag{3.10}
|\{Mf > 2\lambda\}| \leq |\{M[\chi_{\{|f| > \lambda\}}f] > \lambda\}| \leq \frac{3^{n}}{\lambda}\int_{\{|f| > \lambda\}}|f(x)|\dd x
\label{eq:3.10}
\end{equation}
成立. 由式~\eqref{eq:3.8} 作变量替换 $\sigma = 2\lambda$, 得
\begin{equation}\tag{3.11}
\|Mf\|_{L^{p}} = \left(p2^{p}\int_{0}^{\infty}\lambda^{p-1}|\{x \in \mathbb{R}^{n}: Mf(x) > 2\lambda\}|\dd\lambda\right)^{1/p}.
\label{eq:3.11}
\end{equation}
将式~\eqref{eq:3.10} 与式~\eqref{eq:3.11} 结合,
\begin{align*}
\int_{\mathbb{R}^{n}}Mf(x)^{p}\dd x &= p2^{p}\int_{0}^{\infty}\left(\int_{\mathbb{R}^{n}}\chi_{\{Mf(x) > 2\lambda\}}(x,\lambda)\dd x\right)\lambda^{p-1}\dd\lambda \\
&\leq p2^{p}\cdot 3^{n}\int_{0}^{\infty}\left(\int_{\{|f| > \lambda\}}|f(x)|\dd x\right)\lambda^{p-2}\dd\lambda.
\end{align*}
使用 Fubini 定理, 仔细计算一维定积分, 得
\begin{align*}
\int_{\mathbb{R}^{n}}Mf(x)^{p}\dd x
&\leq p2^{p}\cdot 3^{n}\int_{0}^{\infty}\left(\int_{\mathbb{R}^{n}}\chi_{\{(x,\lambda): |f(x)| > \lambda\}}(x,\lambda)|f(x)|\dd x\right)\lambda^{p-2}\dd\lambda \\
&= p2^{p}\cdot 3^{n}\int_{\mathbb{R}^{n}}\left(\int_{0}^{\infty}\chi_{\{(x,\lambda): |f(x)| > \lambda\}}(x,\lambda)\lambda^{p-2}\dd\lambda\right)|f(x)|\dd x \\
&= p2^{p}\cdot 3^{n}\int_{\mathbb{R}^{n}}\left(\int_{0}^{|f(x)|}\lambda^{p-2}\dd\lambda\right)|f(x)|\dd x \\
&= \frac{p2^{p}\cdot 3^{n}}{p-1}\int_{\mathbb{R}^{n}}|f(x)|^{p}\dd x.
\end{align*}
\end{proof}

前面说过, Hardy--Littlewood 极大算子对研究积分算子是有效的, 作为其一例, 来考虑分数次积分算子. 设 $0 < \alpha < n$. 对适当的可测函数 $f$, 定义
\begin{equation}\tag{3.12}
I_{\alpha}f(x) \equiv \int_{\mathbb{R}^{n}}\frac{f(y)}{|x-y|^{n-\alpha}}\dd y \quad (x \in \mathbb{R}^{n}).
\label{eq:3.12}
\end{equation}
由于积分必须被定义, 将 $f$ 取为非负值函数是一种想法, 但可知对 $L^{p}$ 函数这个积分算子也能定义. 下面证明下面的定理.

\begin{theorem}[Hardy--Littlewood--Sobolev]\label{thm:3.5}
当 $1 < p < q < \infty$, $0 < \alpha < n$ 满足 $\frac{1}{q} = \frac{1}{p} - \frac{\alpha}{n}$ 时, 对非负值 Lebesgue 可测函数 $f$, 有 $\|I_{\alpha}f\|_{L^{q}} \lesssim \|f\|_{L^{p}}$ 成立.
\end{theorem}

\begin{proof}
将积分核改写为
\[
\frac{1}{|x-y|^{n-\alpha}} = (n-\alpha)\int_{0}^{\infty}\chi_{B(x,r)}(y)\frac{\dd r}{r^{n-\alpha+1}}.
\]
于是, 由 Fubini 定理, 可以表为
\[
I_{\alpha}f(x) = (n-\alpha)\int_{0}^{\infty}\left(\frac{1}{r^{n-\alpha+1}}\int_{B(x,r)}f(y)\dd y\right)\dd r.
\]
这里, 由 Hardy--Littlewood 极大算子的定义与 Hölder 不等式, 得
\[
\int_{B(x,r)}f(y)\dd y \leq |B(x,r)|Mf(x), \quad \int_{B(x,r)}f(y)\dd y \leq |B(x,r)|^{1/p'}\|f\|_{L^{p}}.
\]
即
\[
\int_{B(x,r)}f(y)\dd y \leq \min(|B(x,r)|Mf(x), |B(x,r)|^{1/p'}\|f\|_{L^{p}}).
\]
这里, 关于正的未知数 $r$ 的方程 $r^{n}A = r^{n/p'}B$ 的解是 $r = \left(\frac{B}{A}\right)^{p/n}$. 所以,
\begin{align*}
\int_{0}^{\infty}\min(r^{n}A, r^{n/p'}B)r^{-n+\alpha-1}\dd r
&= A\int_{0}^{(B/A)^{p/n}}r^{\alpha-1}\dd r + B\int_{(B/A)^{p/n}}^{\infty}r^{-n/p+\alpha-1}\dd r \\
&= A\int_{0}^{(B/A)^{p/n}}r^{\alpha-1}\dd r + B\int_{(B/A)^{p/n}}^{\infty}r^{-n/q-1}\dd r \\
&= \frac{A}{\alpha}\left(\frac{B}{A}\right)^{\frac{p}{n}\alpha} + \frac{nB}{q}\left(\frac{A}{B}\right)^{p/q} \\
&= \left(\frac{1}{\alpha} + \frac{n}{q}\right)B\left(\frac{A}{B}\right)^{p/q} = \left(\frac{1}{\alpha} + \frac{n}{q}\right)A^{p/q}B^{1-p/q}.
\end{align*}
所以
\[
I_{\alpha}f(x) \lesssim Mf(x)^{p/q}\|f\|_{L^{p}}^{1-p/q}
\]
成立, 对此式两边取 $L^{q}$ 范数,
\begin{equation}
\begin{aligned}
\|I_{\alpha}f\|_{L^{q}} &\lesssim \|(Mf)^{p/q}\|_{L^{q}}\|f\|_{L^{p}}^{1-p/q} \\
&= (\|Mf\|_{L^{p}})^{p/q}\|f\|_{L^{p}}^{1-p/q}
\end{aligned}\tag{3.13}
\label{eq:3.13}
\end{equation}
成立. 由于 $1 < p < \infty$, $M$ 是 $L^{p}$ 有界的, $\|Mf\|_{L^{p}} \lesssim \|f\|_{L^{p}}$ 成立. 所以, 将此不等式代入式~\eqref{eq:3.13},
\[
\|I_{\alpha}f\|_{L^{q}} \lesssim (\|f\|_{L^{p}})^{p/q}\|f\|_{L^{p}}^{1-p/q} = \|f\|_{L^{p}}.
\]
通过分解 $f = \max(f, 0) - \max(-f, 0)$, 可得 $I_{\alpha}$ 的 $L^{p}$ 有界性.
\end{proof}

\begin{corollary}[Hardy--Littlewood--Sobolev]\label{cor:3.1}
当 $1 < p < q < \infty$, $0 < \alpha < n$ 满足 $\frac{1}{q} = \frac{1}{p} - \frac{\alpha}{n}$ 时, 对 $f \in L^{p}(\mathbb{R}^{n})$, $I_{\alpha}[|f|](x) < \infty$ 几乎处处成立, 且 $\|I_{\alpha}f\|_{L^{q}} \lesssim \|f\|_{L^{p}}$ 成立.
\end{corollary}

\begin{proof}
$I_{\alpha}[|f|](x) < \infty$ 几乎处处成立, 是由 $I_{\alpha}[|f|] \in L^{q}(\mathbb{R}^{n})$ 得到的. 不等式如上面所述, 通过分解 $f = \max(f, 0) - \max(-f, 0)$ 得到.
\end{proof}

\begin{exercise}\label{exer:3.6}
对 Hardy--Littlewood 极大算子 $M$ 与可测函数 $f$, $g$ 及 $a \in \mathbb{C}$, 证明下面的函数不等式或等式.

(1) $M[f+g] \leq Mf + Mg$. 不过, 可以承认积分的三角不等式.

(2) $M[af] = |a|Mf$.

【注意】应当证明的是, 对任意 $x \in \mathbb{R}^{n}$, (1) $M[f+g](x) \leq Mf(x) + Mg(x)$, (2) $M[af](x) = |a|Mf(x)$ 成立.
\end{exercise}
